\documentclass[12pt]{article}
\usepackage{amsmath}
\usepackage{amssymb}
\usepackage{amsthm}

\newtheorem{theorem}{Theorem}
\newtheorem{remark}{Remark}

\title{Non-negative Supermartingale Theorem}
\author{Short Notes by Iman Taghaddosinejad}


\begin{document}
\maketitle
Non-negative supermartingales arise naturally in dynamic macroeconomic models through 
Euler inequalities and provide a useful way to establish the long-run convergence of 
endogenous stochastic processes.

\begin{theorem}[Non-negative supermartingale convergence theorem]\label{thm:nonnegsupermartingale}
Let $X_t$ be a non-negative supermartingale such that $X_t\geq\mathbb{E}_t[X_{t+1}]$
and $X_t\geq0$ for all $t\geq0$. Then, there exists a finite random variable 
$X_\infty\geq0$ such that $X_t\to X_\infty$ a.s. as $t\to\infty$. Moreover, we have 
$\mathbb{E}[X_\infty]\leq \mathbb{E}[X_0]$.
\end{theorem}

\begin{remark}[Intuition]
Theorem \ref{thm:nonnegsupermartingale} states that if a random process is a 
non-negative supermartingale then it converges almost surely (along any sequence with positive
probability) to a finite non-negative (random variable) limit. The process will not diverge 
to $\pm\infty$ or oscillate over some positive range $b-a$ with $b>a>0$ infinitely 
often. 
\end{remark}
The proof of Theorem \ref{thm:nonnegsupermartingale} intuitively works in three steps. 
First, we rule out divergence to $\pm\infty$ using basic properties of the process itself. 
Second, we rule out persistent oscillations by showing that the process cannot upcross any 
interval of positive width infinitely often. The key observation is that if
$\liminf_{t\to\infty}X_t<\limsup_{t\to\infty}X_t$, then there must exist some interval 
$[a,b]$ of positive width lying strictly between them that the process upcrosses infinitely often. 
Thus, ruling out infinitely many upcrossings rules out any positive gap between the 
$\liminf$ and $\limsup$. Finally, having ruled out both divergence and persistent oscillation, 
we obtain $\liminf_{t\to\infty}X_t=\limsup_{t\to\infty}X_t <\infty$ a.s. and hence 
$X_t$ converges to a finite limit almost surely.

\begin{proof}[\textbf{Proof} (Sketch proof of Theorem \ref{thm:nonnegsupermartingale})]
Let $X_t$ be a non-negative supermartingale. First we establish the process cannot diverge to 
$\pm\infty$. Since $X_t\geq0$ for all $t\geq0$, we immediately have $\liminf_{t\to\infty} 
X_t \geq 0$ a.s. Moreover, since $X_t\geq\mathbb{E}_t[X_{t+1}]$, taking the unconditional 
expectation of both sides and using the law of iterated expectations (LIE) yields 
$\mathbb{E}[X_t]\geq\mathbb{E}[X_{t+1}]$.
Therefore, $\infty>\mathbb{E}[X_0]\geq\mathbb{E}[X_1]\geq\mathbb{E}[X_2]\geq\cdots\geq0$.
Hence, $\mathbb{E}[X_t]$ is a decreasing sequence bounded below by zero 
and therefore converges to some finite limit. Moreover, the sequence is uniformly bounded
above, $\sup_{t\geq0}\mathbb{E}[X_t]=\mathbb{E}[X_0]<\infty$. Together with non-negativity, 
this rules out the event $\mathcal{A}:=\{\omega:X_t(\omega)\to+\infty\}$, having positive
probability. Otherwise, paths in $\mathcal{A}$ would make an arbitrarily large positive 
contribution to the ensemble mean, while $X_t\geq0$ prevents this contribution from being 
offset by negative realizations, contradicting the finite uniform bound on $\mathbb{E}[X_t]$.
Hence, $\mathbb{P}(X_t\to+\infty)=0$. 

Next we rule out persistent oscillations over any interval of positive width. Let $a<b$ 
and suppose that $0\leq\liminf_{t\to\infty}X_t<a<b<\limsup_{t\to\infty}X_t$. Let 
$U_T$ denote the number of completed upcrossings of the interval $[a,b]$ by time $T$ where 
$0\leq U_1\leq U_2\leq\cdots\leq U_T$. 
Further define the auxiliary process $H_t\in\{0,1\}$ such that $H_t=1$ if $X_t$ is in the 
midst of an upcrossing of $[a,b]$ and $H_t=0$ otherwise. At any time $T\geq0$, 
the process has therefore completed $U_T$ upcrossings and may additionally be in the 
midst of at most one unfinished upcrossing. Let $s_i$ and $\tau_i$ denote the start and 
completion times of the $i$-th completed upcrossing, respectively, so that $X_{s_i}\leq a$, 
$X_{\tau_i}\geq b$, and hence $X_{\tau_i}-X_{s_i}\geq b-a$. If an unfinished upcrossing 
remains at time $T$, let $s$ denote its starting time, so that $X_s\leq a$. Then
we can define the sum of all contributions of upcrossings 
{\footnotesize
\begin{align*}
    G_T&=\sum_{t=0}^{T-1}H_t(X_{t+1}-X_t) \\
    &= \underbrace{H_0(X_1-X_0)+\cdots+H_{s-1}(X_s-X_{s-1})}_{\text{completed upcrossings}} 
        + \underbrace{H_s(X_{s+1}-X_s)+\cdots+H_{T-1}(X_T-X_{T-1})}_{\text{unfinished crossing}} \\
    &= \underbrace{\sum_{i=1}^{U_T} (X_{\tau_i}-X_{s_i})}_{\geq (b-a)U_T\geq0} 
        + \underbrace{X_{T}-X_{s}}_{\geq-a} \\
    &\geq (b-a)U_T - a,
\end{align*}
}%
where the last inequality follows because $X_T\geq0$ and $X_s\leq a$. Furthermore, we can derive 
an upper bound on the ensamble mean of $G_T$: 
\begin{align*}
    \mathbb{E}[G_T] &= \sum_{t=0}^{T-1} \mathbb{E}\bigl[H_t(X_{t+1}-X_t)\bigr] \\
    &= \sum_{t=0}^{T-1}\mathbb{E}\Bigl[\mathbb{E}_t\bigl(H_t(X_{t+1}-X_t)\bigr)\Bigr] \\
    &= \sum_{t=0}^{T-1}\mathbb{E}
        \Bigl[
        \underbrace{H_t}_{\geq0} 
        \underbrace{\mathbb{E}_t(X_{t+1}-X_t)}_{\leq 0}
        \Bigr] \\
    &\leq 0,
\end{align*}
where the second equality uses the LIE, the third equality uses the fact that $H_t$ is 
known at time $t$, and the last inequality uses the supermartingale property. 
Combining the pathwise inequality for $G_T$ with the ensemble mean inequality yields
$\mathbb{E}[U_T] \leq \frac{a}{b-a} < \infty$ for all $T\geq0$. Therefore, 
$\sup_{t\geq0}\mathbb{E}[U_t]\leq \frac{a}{b-a} < \infty$.
Since $U_T\geq0$, this uniform bound rules out any positive-probability set of paths 
along which the number of completed upcrossings grows without bound. Thus, for the event
$\mathcal{B}:=\{\omega:U_T(\omega)\to\infty\}$, we have $\mathbb{P}(\mathcal{B})=0$.
Hence, with probability one, $X_t$ can complete only finitely many upcrossings of any 
fixed interval $[a,b]$ of positive width. This contradicts the assumption 
that $\liminf_{t\to\infty}X_t<\limsup_{t\to\infty}X_t$ since any positive gap contains 
an interval of positive width that must be upcrossed infinitely often. Therefore,
$0\leq\liminf_{t\to\infty}X_t=\limsup_{t\to\infty}X_t<+\infty$ a.s., and hence 
$X_{t}\to X_{\infty}<\infty$ a.s. as $t\to\infty$. 
\end{proof}

\begin{remark}[Null sets, countability, and rational intervals]
A subtle point in the upcrossing argument is that we do not establish that the set of paths
which upcross a fixed interval $[a,b]$ infinitely often is empty. Rather, for each
fixed $a<b$, we establish the weaker result that this set has probability zero.
This distinction matters, as we do not rule out the existence of probability-zero paths.

To pass from the result for each fixed interval to almost-sure convergence, let
$\mathcal{C}:=\left\{\omega:\liminf_{t\to\infty}X_t(\omega) <\limsup_{t\to\infty}X_t(\omega)\right\}$, 
and define the event 
$\mathcal{B}_{a,b}:=\left\{\omega:U_T(\omega)\to\infty, \; \text{with} \; U_T \;\text{on}\; [a,b], \; a,b\in\mathbb{Q} \right\}$.
For every path $\omega\in \mathcal{C}$, $\exists$ rationals $a<b$ satisfying 
$\liminf_{t\to\infty}X_t<a<b<\limsup_{t\to\infty}X_t$, and hence $\omega\in\mathcal{B}_{a,b}$,
Therefore,
\[
\mathcal{C}\subseteq
\bigcup_{\substack{a,b\in\mathbb{Q},\\a<b}}\mathcal{B}_{a,b}.
\]
Moreover, $\mathbb{Q}$ is countably infinite, such that the collection of rational pairs $(a,b)$ is also countable. 
Therefore,
\[
\mathbb{P}(\mathcal{C}) 
\leq \mathbb{P}\left(\bigcup_{\substack{a,b\in\mathbb{Q},\\a<b}}\mathcal{B}_{a,b}\right)
\leq\sum_{\substack{a,b\in\mathbb{Q},\\a<b}}\mathbb{P}(\mathcal{B}_{a,b})=0,
\]
where the first inequality uses monotonicity of probability, while the second uses
countable subadditivity since the union is countable. 
Thus, exceptional paths may exist, but the set containing all of them has probability zero.
The use of $\mathbb{Q}$ ensures that every positive gap formed by a path in $\mathcal{C}$ 
contains a rational interval, while being countable enough that the union of the corresponding 
probability-zero events remains a probability-zero event.
\end{remark}



\end{document}